\chapter*{Wstęp}
\addcontentsline{toc}{chapter}{Wstęp}

Rozwój technologii Internetu Rzeczy (IoT) umożliwia zdalne monitorowanie parametrów środowiskowych w czasie rzeczywistym. Jednym z najczęściej mierzonych parametrów jest temperatura, wykorzystywana zarówno w systemach domowych, jak i przemysłowych. Niniejszy projekt przedstawia implementację systemu składającego się z mikrokontrolera ESP32 WROOM z czujnikiem temperatury DHT11, lokalnej bazy danych MySQL (XAMPP) oraz aplikacji mobilnej na system Android.

System umożliwia zapisywanie pomiarów temperatury do bazy danych oraz ich odczyt i prezentację w aplikacji mobilnej w postaci aktualnej wartości oraz dziennika historycznego.
